% ==============================================================================
% CONCLUSION GÉNÉRALE
% Mémoire MAE -- Arthur SAUNIER -- Wavestone Luxembourg
% Le \chapter{} est déclaré dans le fichier maître ; ce fichier démarre en \section.
% ==============================================================================

\section{Le chemin parcouru}

Ce mémoire est né d'une friction observée avant d'être formulée~: dans une practice
dont l'activité repose sur la production intensive d'écrits opposables, des outils capables
de rédiger vite sont apparus sans que rien ne change dans la manière de contrôler ce qui est
rédigé. La question centrale en est sortie --- comment l'automatisation des tâches d'analyse
redéfinit-elle la légitimité et le développement des compétences des consultants~?

Le traitement de cette question a suivi un ordre délibérément inversé. Le terrain a été
décrit d'abord, avec ce qui le contraint~: un marché où la conformité impose un calendrier
exogène, un format écrit et une exigence de justesse asymétrique. Les cadres conceptuels ont
été mobilisés ensuite, non pour illustrer le terrain mais pour le rendre lisible~: la firme
intensive en connaissances et son modèle de transmission, les registres de la légitimité
professionnelle, les travaux empiriques consacrés aux effets de l'usage. Une enquête
qualitative par entretiens semi-directifs, conduite auprès de collaborateurs occupant les
deux positions de la chaîne de production du livrable --- celle qui rédige et celle qui
contrôle~---, a fourni le matériau. Les résultats qui suivent en sont issus.

\section{Ce que le terrain établit}

Sur l'appropriation, le panel converge et converge sur davantage que le principe. Les tâches
déléguées relèvent toutes du traitement d'un matériau déjà constitué~; l'analyse elle-même ne
l'est jamais. Chacun s'est constitué une zone qu'il ne délègue pas, et cette zone se définit
partout de la même façon~: par la sensibilité des données et par la nécessité d'un contexte
que l'outil n'a pas connu, non par la difficulté de la tâche. Chacun a également repéré le
seuil au-delà duquel l'outil coûte plus de temps qu'il n'en fait gagner. Ce que la
littérature nomme frontière dentelée \cite{dellacqua2026} a donc été localisé sur le terrain,
mais par une méthode unique et coûteuse~: l'échec, mesuré en minutes perdues, puis mémorisé
en règle personnelle. Chacun a payé ce prix séparément, sans savoir que ses collègues
l'avaient payé aussi. L'appropriation est réelle~; elle est privée. Le gain de temps, enfin,
ne sort pas de la boucle de production~: il se déplace de la rédaction vers la vérification.
Un répondant fait exception et réinvestit ce temps dans la relation client.

Sur les compétences, le résultat le plus solide est celui qui retourne l'hypothèse de départ.
L'appauvrissement de la relation entre junior et encadrant avait été inscrit dans la grille
d'analyse, en cohérence avec ce qu'établissent les travaux sur les analystes juniors exposés
à une technologie opaque \cite{anthony2021}~: si la machine produit le premier jet, la boucle
de correction perd son objet. Ce n'est pas ce que le terrain donne à voir, et les deux
positions le contredisent indépendamment l'une de l'autre. La relation ne s'appauvrit pas~:
elle change d'objet. La forme étant déléguée, la revue se reporte sur ce qu'elle n'atteignait
auparavant qu'après plusieurs itérations --- la justification, la pertinence, la source. Le
support de la socialisation \cite{nonaka1997} n'a pas disparu, il a changé de niveau, et la
formation à la vérification critique devient une composante du rôle d'encadrement. Ce
déplacement a une contrepartie, que le corpus documente aussi~: il suppose que le junior ait
exécuté au moins une fois manuellement les tâches qu'il délègue ensuite. Sans cette
exécution, il valide sans pouvoir juger.

Sur la légitimité, le matériau autorise moins. La non-déclaration de l'usage au client
converge, avec une justification commune --- le client achète un résultat et non un
processus~---, et le panel situe la valeur future du cabinet du côté du jugement et de la
responsabilité assumée \cite{suchman1995}. Mais cette convergence a été obtenue sur
une question projective, dont la forme appelle une réponse valorisante, et une seule
interpellation réelle par un client est documentée. Ce qui est établi est ce que les
professionnels estiment devoir vendre~; comment la légitimité se reconstruit effectivement
devant un client ne l'est pas.

Sur la gouvernance, le résultat est net et il est de premier ordre. Le cadre en vigueur est
un cadre d'autorisation --- tel outil est permis, tel autre non --- et non un cadre de
vérification~: il délimite des outils, il ne délimite pas des pratiques. Aucune étape de
contrôle propre aux contenus produits avec assistance n'existe. Là où la pratique de revue a
malgré tout évolué, elle l'a fait après un incident vécu, individuellement et sans que
l'adaptation soit écrite, partagée ni opposable~; là où aucun incident marquant n'a été
rencontré, la revue est demeurée inchangée. Les quasi-incidents, eux, ne remontent pas. Le contrôle ne
s'adapte donc pas par la règle, il s'adapte par l'incident. Le circuit de revue, hérité d'une
logique où le contrôle porte sur la personne formée plutôt que sur l'acte
\cite{mintzberg1982}, tient tant que le relecteur sait quoi chercher~; or la nature des
erreurs a changé sans que le contenu de la qualification ait suivi \cite{lebovitz2022}. Un
dispositif dont le seul mécanisme d'apprentissage est l'accident est un dispositif qui attend
l'accident.

\section{Ce que le travail apporte au cabinet}

Les recommandations formulées ne procèdent pas d'une opinion sur l'opportunité de ces outils.
Elles reprennent ce que le terrain a lui-même énoncé, en le rendant opérationnel, et se
tiennent au point d'équilibre sur lequel le panel converge~: ni interdiction, ni
laisser-faire. Trois lignes s'en dégagent. La première rend la vérification déclarative
plutôt que détective, en graduant la relecture selon ce que le producteur signale --- il ne
s'agit pas de rechercher l'usage, mais de permettre de le déclarer sans coût. La deuxième
déplace l'effort de formation de la formulation des consignes vers la détection des erreurs,
et l'étend aux encadrants, dont la compétence de relecture est celle que le changement de
nature des erreurs a rendue obsolète. La troisième conditionne les deux autres~: le
dispositif doit apprendre autrement que par l'accident, ce qui suppose un traitement organisé
des quasi-incidents, dissocié de toute conséquence individuelle. Un incident détecté avant le
client est une information de valeur~; sa remontée doit être plus avantageuse que son
silence.

\section{La portée des résultats}

Ces résultats sont indexés à leur terrain. Le panel est resserré, il porte sur une seule
practice, d'une seule entité, sur un seul marché, dont la structure --- place
financière concentrée, pression réglementaire spécifique --- n'est pas transposable telle
quelle. Le matériau est constitué de prises de notes et non de transcriptions~: les extraits
reproduits n'ont pas valeur de citation littérale. Les convergences relevées ne valent donc
pas comme fréquences, et rien ici ne prétend décrire la profession du conseil ni même
l'ensemble du cabinet. La portée de l'ensemble relève de la transférabilité analytique vers
des organisations comparables, non de la généralisation \cite{lincoln1985}. Ce que ce travail
propose est un mécanisme, offert à l'examen d'organisations qui reconnaîtraient leur propre
situation dans celle qui est décrite.

\section{Une question que ce travail rend visible sans la trancher}

Un dernier point mérite d'être versé au dossier, précisément parce que l'enquête ne l'a pas
traité. Interrogés sur ce qui n'avait pas été abordé, des répondants ont désigné la même zone
--- la dimension juridique de ces usages~: propriété intellectuelle des contenus produits,
confidentialité des données confiées à un tiers, et surtout imputation de la responsabilité
lorsqu'une erreur non détectée parvient au client. Le constat qui l'accompagne est celui d'un
angle mort~: personne n'en parle vraiment, et le sujet est appelé à devenir central.

L'observation a une portée qui dépasse l'anecdote. Toutes les recommandations formulées ici
reposent sur une même hypothèse implicite~: que le consultant qui signe demeure celui qui
répond. C'est ce que le terrain affirme sans hésiter, et c'est ce qui fonde la traçabilité
déclarative comme la formation à la vérification. Mais cette imputation est un fait
juridique autant qu'un principe professionnel, et rien dans ce mémoire ne l'établit. Si elle
devait se révéler moins évidente qu'il n'y paraît --- entre le cabinet, son collaborateur, le
fournisseur de l'outil et le client~---, l'architecture managériale proposée resterait
pertinente, mais elle ne serait plus suffisante. Cette question relève d'un autre travail,
conduit avec d'autres compétences que celles mobilisées ici. Elle mérite de l'être avant que
ce ne soit un incident qui l'impose.